\documentclass[10pt,conference]{IEEEtran}
\usepackage[T1]{fontenc}
\usepackage{lmodern}
\usepackage{microtype}
\usepackage{amsmath,amssymb}
\usepackage{booktabs}
\usepackage{tabularx}
\usepackage{array}
\usepackage{graphicx}
\usepackage{enumitem}
\usepackage[hidelinks]{hyperref}
\setlist{nosep,leftmargin=*}
\newcolumntype{Y}{>{\raggedright\arraybackslash}X}
\newcommand{\DCCEM}{DCCEM}

\title{Why Are Most People Right-Handed?\\
Developmental Canalization and Context Expression Model}
\author{Source-Controlled Research Plan\\
Survey--Idea--ExperimentDesign--Author Pipeline\\
Developmental neuroscience and quantitative genetics}

\begin{document}
\maketitle

\begin{abstract}
Across populations, most people prefer the right hand for skilled manual tasks, yet the mechanism producing this majority is not a simple Mendelian trait. This paper converts the broad question ``Why are most people right-handed?'' into a falsifiable developmental systems problem. The proposed Developmental Canalization and Context Expression Model (DCCEM) separates a latent neurodevelopmental lateralization state from its observable expressions: hand preference, bilateral motor-performance asymmetry, and writing-hand choice. It integrates common polygenic liability, rare coding variation, developmental covariates, multivariate brain asymmetry, and measured task or cultural exposure. The Survey stage establishes that left-handedness is approximately 10.6\% under a broad meta-analytic estimate, that prevalence is sensitive to measurement and participant characteristics, that large studies implicate many central-nervous-system and microtubule-related loci, and that forced switching can alter estimates [1]--[6]. The Idea stage rejects a single gene, a single cortical location, and a culture-only account, selecting DCCEM as a stronger but testable synthesis. ExperimentDesign specifies a family-aware longitudinal program using repeated multi-task phenotyping, genotyping, optional exome data, non-invasive MRI, and exposure measurements. The primary tests ask whether early latent state predicts later phenotype out of sample, whether brain asymmetry mediates genetic liability, and whether culture changes preference more than performance. This is a design-only proposal: it reports no newly collected data, does not provide an individual diagnostic score, and does not claim that right-handedness is biologically superior.
\end{abstract}

\begin{IEEEkeywords}
handedness, laterality, developmental neuroscience, polygenic architecture, brain asymmetry, gene--environment interaction, motor control
\end{IEEEkeywords}

\section{Introduction}
Handedness is often presented as a binary fact: roughly nine out of ten people are right-handed. The ratio is scientifically interesting precisely because it is neither 50:50 nor universal. A majority is stable enough to appear across time and place, while a sizeable minority is left-handed, mixed-handed, or task-dependent. The scientific problem is therefore not merely to find a correlate of writing hand. It is to explain how a population-level bias emerges from a developing nervous system while preserving individual and task-level variation.

The question also contains a measurement trap. Preference is the hand a person chooses; performance is the relative speed, accuracy, or coordination of each hand; writing hand is a socially visible and historically trained behavior. These variables can agree, but they need not. A person who was encouraged to write with the right hand can still have a different motor-performance asymmetry, and a person with weak preference can select a hand according to tool layout or task demands. Treating all three as one binary phenotype can make genetic or cultural explanations appear stronger than they are.

This paper proposes the Developmental Canalization and Context Expression Model (DCCEM). Its core claim is that a distributed developmental system creates a probabilistic population bias toward right-hand preference, while polygenic and rare-variant liability, early developmental conditions, hemispheric organization, and task environment determine how that bias is expressed. The claim is deliberately stronger than a list of limitations: it predicts a measurable latent state, a temporal relation among state, brain, performance, and preference, and a specific asymmetry between cultural effects on reported preference and effects on bilateral performance.

The proposal has four objectives. First, it compiles evidence into a traceable gap ledger rather than treating the prompt's 85--90\% estimate as a mechanism. Second, it compares alternative Idea routes and records why DCCEM is selected. Third, it formalizes a longitudinal, family-aware design that can distinguish genetic liability from learned expression. Fourth, it defines outcomes that would support, weaken, or falsify the model. The document remains a research plan, not an empirical article and not a clinical or educational decision tool.

\section{Background, Survey, and Research Gap}
\subsection{Prevalence is robust but operationally conditional}
Papadatou-Pastou \textit{et al.} synthesized five meta-analyses and estimated left-handedness at 10.6\% overall. Their estimates ranged from 9.3\% under a stringent left-handedness criterion to 18.1\% under the more permissive non-right-handedness criterion [1]. Study year, assessment method, sex, and ancestry moderated the estimate. The review also identified cultural pressure against left-hand use as a moderator, while arguing that the same broad population ratio persists across regions [1]. Thus, the right-hand majority is a dependable population pattern, but its measured magnitude is partly a property of the instrument and social history.

The implication for DCCEM is immediate: a primary endpoint cannot be ``writing hand'' alone. A valid phenotype must retain item-level responses, task familiarity, spontaneous choice, and performance asymmetry. It must allow right, left, mixed, and weak-preference classes rather than forcing ambiguous observations into a binary label.

\subsection{Family studies support probability, not destiny}
In a nationwide Finnish sample of 30,161 adults, Vuoksimaa \textit{et al.} examined handedness, forced right-handedness, and twin status [5]. Their model estimates attributed most observed variance to unique environmental effects, with estimates dependent on developmental stage and model specification. Forced right-handedness changed estimates of familial effects. This result is not evidence that genes are irrelevant; it is evidence that phenotype history and social correction must be measured before interpreting familial resemblance.

The twin meta-analysis by Pfeifer \textit{et al.} included 59 studies. It found higher raw prevalence of left- and non-right-handedness among twins than singletons, but no meaningful monozygotic-versus-dizygotic prevalence difference. Monozygotic concordance was only modestly elevated, and the authors concluded that secular trends and covariates make a genuine twinning effect unsafe to infer [4]. A recent genetics review summarizes handedness heritability at approximately 25\% [6]. These findings are compatible with a distributed genetic liability that is shaped by development and expression context; they are incompatible with a deterministic ``right-handedness gene.''

\subsection{Genome-wide and rare-variant evidence}
The largest GWAS synthesis in the Survey combined UK Biobank, 23andMe, and the International Handedness Consortium, totaling 1,766,671 participants. It identified 41 loci associated with left-handedness and 7 associated with ambidexterity [2]. Tissue enrichment implicated the central nervous system, while pathway analyses highlighted microtubule regulation and brain morphology. Importantly, left-handedness and ambidexterity showed a low genetic correlation, indicating that they should not be treated as opposite ends of one simple scale.

Wiberg \textit{et al.} related handedness to brain imaging phenotypes in approximately 9,000 UK Biobank participants and performed GWAS in approximately 400,000 participants [3]. Their significant loci included genes or expression quantitative trait loci related to brain development and patterning, including MAP2, MAPT, WNT3, and MICB. They reported altered connectivity between left and right language networks in left-handers. These results identify a plausible developmental bridge but do not establish that language causes handedness or that one tract contains handedness.

Schijven \textit{et al.} extended the architecture to rare coding variation by analyzing exome data from 38,043 left-handers and 313,271 right-handers [7]. Rare coding variants in TUBB4B were 2.7 times more frequent in left-handers, with additional evidence for DSCAM and FOXP1; the estimated exome-wide heritability from rare coding variants was 0.91\% [7]. The finding strengthens the microtubule and neurodevelopmental hypothesis while quantifying its limited explanatory share. DCCEM therefore includes rare variants as one liability layer, not as a new deterministic label.

\subsection{Brain asymmetry as a candidate mediator}
Brain imaging studies provide an intermediate level between DNA and behavior. The results above associate handedness and genetic influences on handedness with structural or functional asymmetries [3], [8], [9]. A 31,864-person study linked handedness and its genetic influences to structural asymmetries of the cerebral cortex [8], and polygenic-score work examined related asymmetries in brain structure [9].

The evidence does not justify a one-region localization. Brain asymmetry is multivariate, developmentally dynamic, and measurement-dependent. DCCEM consequently treats $A_{it}$ as a vector of candidate mediator features whose reliability must be established before mediation is attempted. It asks whether a stable combination of motor, language-related, and interhemispheric features carries information from genetic liability to later performance, while allowing that different individuals can reach similar hand preferences through different neural configurations.

\subsection{Accepted research gaps}
The Survey accepts seven downstream gaps: \textbf{G1}, separation of preference, performance, writing hand, mixed use, and ambidexterity; \textbf{G2}, the developmental mechanism connecting many loci to lateralization; \textbf{G3}, the functional meaning of rare coding contributions; \textbf{G4}, a reliable multiscale brain mediator; \textbf{G5}, gene--environment interaction and forced switching; \textbf{G6}, longitudinal directionality; and \textbf{G7}, an integrated explanation of majority stability with minority persistence. The evidence supports these as open problems, not as hidden findings.

\begin{table}[t]
\caption{Evidence anchors and inference boundaries.}
\label{tab:evidence}
\centering
\small
\begin{tabularx}{\columnwidth}{@{}p{0.07\columnwidth}p{0.27\columnwidth}YY@{}}
\toprule
ID & Anchor & Supported contribution & Boundary carried forward \\
\midrule
E1 & Prevalence meta-analysis [1] & About 10.6\% left-handedness; moderators & No universal rate or phenotype \\
E2 & Large GWAS [2] & 41 left and 7 ambidexterity loci & Polygenic association, not causation \\
E3 & Imaging genetics [3] & Brain-development loci and network differences & No single region or language cause \\
E4 & Exome study [7] & TUBB4B and small rare-variant component & No deterministic coding explanation \\
E5 & Twin meta-analysis [4] & Modest concordance and covariate sensitivity & Heritability is not mechanism \\
E6 & Finnish cohort [5] & Forced switching changes estimates & Environment is not the whole cause \\
E7 & Neurogenetic review [10] & Developmental and evolutionary synthesis & Evolutionary mechanism unsettled \\
E8--E9 & Structural and polygenic imaging [8], [9] & Candidate multivariate mediator & Mediation requires temporal validation \\
\bottomrule
\end{tabularx}
\end{table}

\section{Research Questions and Planned Contributions}
The primary research question is: can a latent, developmentally anchored lateralization state explain why right-hand preference is the population majority while distinguishing the effects of inherited liability, brain organization, and task environment? Four planned contributions follow.

First, the project defines a multitask phenotype rather than a single label. Second, it estimates a path from polygenic and rare-variant liability through brain asymmetry to later motor phenotype. Third, it quantifies context expression by comparing reported preference, writing hand, and bilateral performance under measured histories of instruction and tool affordance. Fourth, it evaluates whether the same model generalizes across an independent cohort and ancestry groups without treating effect-size equality as a requirement.

The contribution is explanatory rather than diagnostic. A model can explain a population distribution and still be unable to predict one person's hand preference with useful accuracy. This distinction is preserved in the analysis plan and in the Author-stage wording.

\section{Idea Source Checkpoints and Direction Selection Audit}
The Idea Agent generated five routes. A single-pathway genetic explanation was rejected because the literature is polygenic and because rare coding variation explains a small share. A single-region cortical localization was rejected because imaging evidence is distributed and heterogeneous. A culture-only explanation was retained as a context module because forced switching and prevalence moderation are real, but it cannot account for neurodevelopmental and family evidence. A frequency-dependent population model was retained as a population module because it addresses majority stability, but it needs individual-level developmental mediators.

The selected route, DCCEM, combines these modules into a hierarchy. Genetic variation affects developmental liability; early developmental conditions and family context shape the trajectory; brain organization mediates part of the trajectory; and task environment shapes expression. The route is not a claim that each edge is already proven. It is a compact arrangement of evidence-aligned hypotheses that can fail at several independent checkpoints.

\begin{table}[t]
\caption{Idea route audit.}
\centering
\small
\resizebox{\columnwidth}{!}{%
\begin{tabular}{lccccc}
\toprule
Route & Novelty & Mechanism & Translation & Gaps & Decision \\
\midrule
Single gene & Low & Low & Low & 2/7 & Reject \\
Single region & Medium & Medium & Medium & 3/7 & Reject \\
Culture only & Medium & Medium & High & 3/7 & Component \\
Evolution ratio & High & Medium & Medium & 3/7 & Component \\
DCCEM & High & High & High & 7/7 & Primary \\
\bottomrule
\end{tabular}}
\end{table}

The four checkpoints are: (C1) binary phenotype to latent multitask phenotype; (C2) locus list to polygenic and rare-variant developmental liability; (C3) adult association to longitudinal genotype--brain--performance--preference ordering; and (C4) culture as a competing explanation to context expression moderation. The primary route is supported only if the locked latent construct is reliable and improves held-out explanation over simpler baselines.

\section{Problem Definition, Assumptions, and Hypotheses}
The unit of analysis is a participant-specific developmental state, not a handedness label alone. Let $L_{it}$ be latent lateralization for participant $i$ at age $t$, $G_i$ a discovery-trained polygenic score, $R_i$ a preregistered rare coding burden, $P_{it}$ developmental covariates, and $u_{f(i)}$ and $v_{s(i)}$ family and site effects:

\begin{equation}
L_{it}=\alpha+\beta_gG_i+\beta_rR_i+\beta_pP_{it}+u_{f(i)}+v_{s(i)}+\epsilon_{it}.
\end{equation}

For task $j$, the observed phenotype $Y_{ijt}$ is modeled as a multiclass outcome with distinct loadings for preference, performance, and writing hand:

\begin{equation}
\begin{aligned}
\Pr(Y_{ijt}=k)=\operatorname{softmax}_k(&\lambda_{jk}L_{it}+\gamma_{jk}E_{ijt}\\
&+\delta_{jk}L_{it}E_{ijt}+c_{ijt}),
\end{aligned}
\end{equation}

where $E_{ijt}$ contains task familiarity, instruction, tool orientation, and measured cultural pressure. The interaction term is central: it tests whether context expression depends on developmental liability.

The brain mediator is a vector $A_{it}$ of structural, diffusion, and functional asymmetry features:

\begin{equation}
A_{it}=\theta_0+\theta_gG_i+\theta_rR_i+\theta_pP_{it}+\theta_eE_{it}+b_i+\zeta_{it}.
\end{equation}

Family data provide an ACE-style sensitivity analysis, $\operatorname{Var}(L)=A+C+E$, with forced-switching and birth covariates. This decomposition is descriptive and model-dependent; it is not treated as a direct intervention target.

The hypotheses are:
\begin{enumerate}
\item H1: a repeated multitask latent phenotype is more reliable and predictive than writing hand alone.
\item H2: polygenic and rare coding terms explain small incremental variance after family and developmental adjustment.
\item H3: multivariate brain asymmetry partially mediates genetic liability without requiring one cortical location.
\item H4: cultural and task exposure shifts preference and writing hand more than bilateral performance asymmetry, with moderation by $L_{it}$.
\item H5: latent developmental stability increases with age while task-specific preference remains more plastic.
\item H6: the direction of these relations generalizes to held-out cohorts and ancestry groups, even when magnitudes differ.
\end{enumerate}

\section{Study Design and Methods}
\subsection{Cohorts and schedule}
The planning design combines existing cohorts with a prospective non-invasive panel. The discovery target is approximately 30,000 participants with genotype data and harmonizable phenotypes; the replication target is approximately 20,000 independent participants. A prospective panel of approximately 1,500 children would be enrolled at 6--18 months and followed at 3, 5, 8, 12, and 16 years. A nested imaging panel of approximately 500 participants would be assessed at 8, 12, and 16 years when feasible. These are planning targets, not observed sample sizes.

Twins and siblings are included where available. Zygosity, birth weight, gestational age, sex, injury, developmental history, and forced-switching history are retained. Their inclusion prevents the common mistake of treating family structure and perinatal factors as unmeasured noise.

\subsection{Phenotyping}
At each age-appropriate visit, preference is recorded across several manual tasks, including reach, object manipulation, drawing, and tool use. Unimanual and bimanual performance is measured with accuracy-matched timing and repeated trials. Writing or drawing hand is recorded separately. Video or inertial-sensor coding is used where feasible to estimate spontaneous choice, switching, movement initiation, and task-specific asymmetry. Footedness and eyedness are secondary laterality measures, not substitutes for handedness.

The measurement model preserves item-level data and distinguishes culturally sensitive tasks from neutral tasks. A latent phenotype is trained only in discovery data, with test--retest reliability assessed before genetic or mediation analyses. Participants are not asked to train or change their preferred hand.

\subsection{Brain and genetic measurements}
The nested imaging panel uses non-invasive structural MRI, diffusion imaging, and task or resting-state functional connectivity. Candidate features include cortical asymmetry, interhemispheric tracts, and motor or language-related network connectivity. Scanner, site, motion, preprocessing batch, and incidental-finding procedures are recorded. No single region is required to enter the mediator; a multivariate feature space is prespecified and dimension-reduced within training data.

Genetic data include genome-wide arrays with imputation and quality control, ancestry-aware principal components, relatedness, and optional exome sequencing. Polygenic scores are trained in discovery data and locked before replication. Rare coding burden is restricted to predeclared allele-frequency and consequence classes, with TUBB4B, DSCAM, and FOXP1 evaluated as registered hypotheses rather than automatic clinical markers. No participant receives an individual genetic handedness score.

\subsection{Context and developmental covariates}
The study records correction or encouragement to use a particular hand, school writing instruction, utensil and workstation orientation, family handedness, language environment, socioeconomic context, birth variables, and major developmental events. These variables are measured as exposures, moderators, or sources of measurement bias; they are not treated as judgments about families or cultures. Cross-cohort harmonization retains the original response scale and includes a missingness indicator when an item cannot be aligned.

\subsection{Analysis and baselines}
The primary analysis locks the phenotype ontology, genetic score construction, and mediator feature rules in discovery data, then evaluates the independent cohort without threshold retuning. Endpoints are multitask reliability, incremental genetic variance, brain mediation under temporal ordering, liability-by-context interaction, age-dependent stability, and transportability. Held-out log loss, Brier score, calibration slope, variance components, and test--retest reliability are primary metrics. AUC is secondary and descriptive because the goal is explanation of a population distribution, not diagnosis.

Competing models include writing-hand-only, environment-only, genetics-only, brain-only, and unconstrained multi-factor baselines. Multiple testing is controlled within prespecified endpoint families. Age, sex, ancestry, site, motion, birth variables, and attrition are modeled explicitly. Missing data are handled with a cohort- and age-aware model, followed by sensitivity analyses for informative dropout. Mediation is reported only when the putative mediator is reliable and precedes the later outcome.

\subsection{Falsification rules}
The DCCEM route is weakened if the latent state fails out-of-sample prediction, if genetic terms fail replication or add no variance, if culture shifts bilateral performance as much as reported preference, if the brain mediator is not reliable across sessions or sites, or if the multi-task phenotype performs no better than a single-task label. A statistical association without temporal identification is not labeled causal. If the model does not generalize, the study returns to the narrower conclusion that current handedness phenotypes do not support a unified developmental explanation.

\subsection{Reproducibility and data partitioning}
The design uses a locked sequence so that the same data cannot be used to invent the phenotype and validate it. First, a data dictionary records the original task, response scale, instruction history, missingness code, and acquisition site. Second, the discovery cohort is split by family, not by individual, into training, tuning, and internal holdout partitions. All transformations that can learn from the data, including imputation parameters, feature selection, dimensionality reduction, polygenic-score thresholds, and latent-factor rotations, are fitted inside the training partition. Tuning choices are frozen before the internal holdout is opened.

The independent replication cohort is then used once for the primary confirmatory analysis. A second external cohort, if available, is reserved for transportability rather than repeatedly used to improve fit. Imaging features are harmonized with site-aware models and their reliability is assessed on repeat scans or split-half data. Genetic analyses retain a separate ancestry-aware calibration report so that a lower prediction score in one group is not mistaken for a biological absence of handedness architecture.

Every reported endpoint is linked to a preregistered estimand, input variables, exclusion rule, and decision threshold. The analysis repository records software versions, quality-control summaries, model seeds, and a machine-readable manifest of derived variables. Where participant privacy prevents releasing genotype or MRI data, synthetic examples and executable analysis specifications are released, while access-controlled data are audited separately. This is necessary because an apparently precise handedness model can otherwise be reproduced only in the cohort in which it was tuned.

The proposed mediation analysis has an additional gate. A brain feature enters the mediator set only when its scan-rescan reliability, missingness pattern, and site calibration are acceptable. The genetic score is trained without access to later brain or outcome measurements. The mediator model is then evaluated against a no-mediator model and a direct-effect model in held-out data. A statistically nonzero indirect path is reported as compatible with mediation, not as proof of a biological mechanism, unless the temporal and identification assumptions are explicitly defended.

This partitioning also protects the majority-stability test. The model must predict the full distribution of right, left, mixed, and weak-preference classes under alternative task definitions. It is not allowed to achieve a good fit by imposing a fixed prior near 90\% right-handedness. Posterior predictive checks therefore vary task composition, cultural exposure, age, and ancestry. A model that fits the aggregate ratio but fails these distributional checks has not explained why the majority persists; it has only restated the input prevalence.

\section{Expected Outcome Branches and Conditional Conclusions}
If H1--H4 replicate, the result would support a distributed developmental account in which the right-hand majority is a population-level bias and culture primarily shapes expression. The appropriate conclusion would still be probabilistic: the model explains how a bias can be generated without predicting every person's hand.

If only the context effect replicates, the proposal retains a preference-expression model but drops the claim that measured brain asymmetry mediates genetic liability. If only genetic associations replicate, the conclusion remains polygenic liability without a demonstrated developmental chain. If brain asymmetry replicates but does not mediate later behavior, it remains a correlated phenotype. If no model generalizes, measurement heterogeneity becomes the principal result and the DCCEM synthesis is rejected.

The population-level majority can be evaluated using posterior predictive distributions. A useful model should reproduce a right-heavy distribution while retaining non-right and mixed classes across task contexts. It should not achieve fit by forcing a fixed 90:10 prior. The model is considered informative only if its predicted distribution changes appropriately when phenotype definition, culture exposure, and developmental stage are varied.

\section{Risks, Limitations, and Human Review Requirements}
The design is ambitious but has identifiable risks. Cohort data may overrepresent particular ancestries, education levels, or health profiles. Handedness instruments may change across ages. MRI motion and site effects can create apparent asymmetry. Polygenic scores can transport poorly across ancestry groups. Rare-variant analyses are sensitive to annotation and coverage. Longitudinal attrition can bias developmental trajectories. These risks require preregistered QC, independent replication, site-aware modeling, and transparent reporting of missingness.

The design is observational and cannot by itself establish all causal edges. A family model cannot perfectly separate genetic and shared environmental effects. Mediation depends on temporal ordering and no-unmeasured-confounding assumptions that may be imperfect. Cultural exposure is not a single variable, and the phrase ``pressure'' must be operationalized without stigmatizing communities. The design therefore demands expert review of the phenotype ontology, ancestry handling, consent and recontact, imaging QC, and interpretation language.

The study is non-invasive, but genetic privacy and child participation require heightened safeguards. Parents and children need age-appropriate consent and assent, with clear policies for withdrawal, recontact, incidental findings, and data deletion. No output should be used to select educational treatment, infer intelligence, or classify a child as neurologically normal or abnormal. Right- and left-handedness are treated as human variation, not a hierarchy.

\section{Conclusion}
Most people are right-handed because a population-level developmental bias is expressed through a heterogeneous system, not because all people carry one right-handedness gene. Existing evidence supports a real majority, polygenic and rare-variant contributions, developmental brain asymmetries, and context-sensitive measurement. It does not yet supply a complete causal chain.

DCCEM turns that unresolved chain into a testable program. Its decisive experiment is not another isolated prevalence survey or another uncontextualized association scan. It is a longitudinal comparison of preference, performance, writing hand, genotype, brain asymmetry, and measured task environment, evaluated in an independent cohort. The proposed direction is bold enough to predict where effects should differ and constrained enough to fail if those predictions do not replicate. This manuscript is a source-controlled proposal with no observed results and no clinical recommendation.

\appendices
\section{Variables and Operational Definitions}
\begin{table}[h]
\caption{Operational variables inherited across the pipeline.}
\centering
\small
\begin{tabularx}{\columnwidth}{@{}p{0.22\columnwidth}YY@{}}
\toprule
Variable & Operational definition & Interpretation rule \\
\midrule
Hand preference & Repeated task-specific choice with item-level responses & Not equivalent to writing hand \\
Performance asymmetry & Accuracy-matched relative speed, force, or coordination across hands & Primary biological expression candidate \\
Writing hand & Hand selected for writing or drawing under recorded instruction history & Socially and task influenced \\
Latent state $L$ & Reliability-checked multitask factor across age & Population model, not diagnosis \\
Brain asymmetry $A$ & Multivariate structural, diffusion, and functional features & Mediator only with temporal validation \\
Genetic liability & Locked polygenic score and rare coding burden & Probabilistic, no individual label \\
Context exposure $E$ & Correction, instruction, tool affordance, and task familiarity & Moderator or measurement source \\
Stability & Cross-age and cross-session latent-state reliability & Not proof of genetic causation \\
\bottomrule
\end{tabularx}
\end{table}

\section{Gap-to-Test Matrix}
The following matrix makes the Survey-to-ExperimentDesign handoff explicit. Each accepted gap has a primary measurement, a model test, and a failure interpretation. This prevents the Author stage from presenting the DCCEM synthesis as if every edge had already been established by the component literature.

\begin{table}[h]
\caption{Accepted gaps mapped to planned tests.}
\centering
\scriptsize
\begin{tabularx}{\columnwidth}{@{}p{0.07\columnwidth}p{0.28\columnwidth}YY@{}}
\toprule
Gap & Measurement & Planned test & Failure interpretation \\
\midrule
G1 & Multi-task preference, performance, and writing hand & Reliability and latent-class comparison & One label is insufficient \\
G2 & Polygenic score and developmental covariates & Incremental variance and replication & No demonstrated genetic path \\
G3 & Rare coding burden and registered genes & Burden and gene-set sensitivity & Rare variants remain minor or unstable \\
G4 & Structural, diffusion, and functional asymmetry & Reliability-gated mediation & Asymmetry is correlated but not mediating \\
G5 & Correction, instruction, tools, and task familiarity & Liability-by-context interaction & Context shifts expression only or not at all \\
G6 & Repeated measures from infancy to adolescence & Cross-age stability and temporal ordering & No developmental canalization evidence \\
G7 & Full class distribution across cohorts & Posterior predictive majority-stability check & Ratio is restated, not explained \\
\bottomrule
\end{tabularx}
\end{table}

The matrix also determines the minimum evidence needed for a positive interpretation. A replicated genetic association without G1 and G6 cannot distinguish a biological liability from a measurement artifact. A reliable brain asymmetry without G3 or temporal ordering cannot establish mediation. A context interaction without bilateral performance measures cannot show whether culture changes motor organization or only the visible task expression. Finally, a model that predicts the right-hand majority but cannot preserve mixed and non-right classes fails the population-level question even if its aggregate fit is high.

\section{Evidence Coverage and Review Appendix}
The Survey-to-Idea binding is recorded in \texttt{survey/survey\_manifest.json}, the selected direction in \texttt{idea/idea\_result.json}, and the design contract in \texttt{experiment\_design/research\_design.json}. The Author stage is constrained to these artifacts. All future empirical outcomes remain an empty list. No acknowledgments section is included. Human review is required before any recruitment, genotyping, imaging, data linkage, or public interpretation.

\begin{thebibliography}{10}
\bibitem{papadatou2020} M. Papadatou-Pastou, E. Ntolka, J. Schmitz, M. Martin, M. R. Munafo, S. Ocklenburg, and S. Paracchini, ``Human handedness: A meta-analysis,'' \textit{Psychological Bulletin}, vol. 146, no. 6, pp. 481--524, 2020, doi: 10.1037/bul0000229.
\bibitem{cuellar2020} G. Cuellar-Partida \textit{et al.}, ``Genome-wide association study identifies 48 common genetic variants associated with handedness,'' \textit{Nature Human Behaviour}, vol. 5, no. 1, pp. 59--70, 2020, doi: 10.1038/s41562-020-00956-y.
\bibitem{wiberg2019} A. Wiberg \textit{et al.}, ``Handedness, language areas and neuropsychiatric diseases: insights from brain imaging and genetics,'' \textit{Brain}, vol. 142, no. 10, pp. 2938--2947, 2019, doi: 10.1093/brain/awz257.
\bibitem{pfeifer2022} L. S. Pfeifer \textit{et al.}, ``Handedness in twins: meta-analyses,'' \textit{BMC Psychology}, vol. 10, art. no. 11, 2022, doi: 10.1186/s40359-021-00695-3.
\bibitem{vuoksimaa2009} E. Vuoksimaa, M. Koskenvuo, R. J. Rose, and J. Kaprio, ``Origins of handedness: a nationwide study of 30,161 adults,'' \textit{Neuropsychologia}, vol. 47, no. 5, pp. 1294--1301, 2009, doi: 10.1016/j.neuropsychologia.2009.01.007.
\bibitem{ocklenburg2025} S. Ocklenburg, A. Mundorf, J. Peterburs, and S. Paracchini, ``Genetics of human handedness: microtubules and beyond,'' \textit{Trends in Genetics}, vol. 41, no. 6, pp. 497--505, 2025, doi: 10.1016/j.tig.2025.01.006.
\bibitem{schijven2024} D. Schijven, S. Soheili-Nezhad, S. E. Fisher, and C. Francks, ``Exome-wide analysis implicates rare protein-altering variants in human handedness,'' \textit{Nature Communications}, vol. 15, art. no. 2632, 2024, doi: 10.1038/s41467-024-46277-w.
\bibitem{sha2021} Z. Sha \textit{et al.}, ``Handedness and its genetic influences are associated with structural asymmetries of the cerebral cortex in 31,864 individuals,'' \textit{Proceedings of the National Academy of Sciences}, vol. 118, no. 47, art. no. e2113095118, 2021.
\bibitem{ocklenburg2021} S. Ocklenburg \textit{et al.}, ``Polygenic scores for handedness and their association with asymmetries in brain structure,'' \textit{Brain Structure and Function}, vol. 227, no. 2, pp. 515--527, 2021.
\bibitem{ocklenburg2013} S. Ocklenburg, C. Beste, and O. Gunturkun, ``Handedness: a neurogenetic shift of perspective,'' \textit{Neuroscience and Biobehavioral Reviews}, vol. 37, no. 10, pp. 2788--2793, 2013.
\end{thebibliography}

\end{document}
